Während des Abbremsens wirken auf den Turner die Normalkraft $F_\text{N} = 2\,m\,g$ nach oben und die Gewichtskraft $m\,g$ nach unten. Die mittlere Bremsbeschleunigung (nach oben) ist damit
\begin{align}
 a &= \frac{F_\text{N}-m\,g}{m} = \aLaF = g = \aLa
\end{align}
Bis die Geschwindigkeit $v$ auf null abgebaut ist, dauert es
\begin{align}
 t &= \frac{v}{a} = \tLaF \\
   &= \frac{\vLa}{\gE} \\
   &= \tLa \approx \result{\tLaP{0}{2}}
\end{align}
Der Abbremsvorgang muss also mindestens so lange dauern. Der Turner verlängert ihn, indem er beim Aufsetzen tief in die Knie geht (und in den Hüften nachgibt), statt mit gestreckten Beinen zu landen. Zusätzlich verlängert eine weiche Landematte die Bremsdauer.
